\documentclass[a4paper]{article}     
\usepackage{amsfonts}
\usepackage{amsmath}
\usepackage{amssymb}
\usepackage{graphicx}

\begin{document}

\noindent \large Auteur: Mateusz Ragan \\
\noindent \large Studierichting: Informatica\\
\noindent \large Oefeningenreeks 2E nr. 26.22\\

\medskip

\normalsize

$\log_{5}(6-x)+\log_{25}x^2\cdot\log_{x}(x+2) = \log_{x}x + \log_{5}x$\\

$\Leftrightarrow \log_{5}(6-x) + \dfrac{2 \log_{5}x}{\log_{5}25} \cdot \dfrac{\log_{5}(x+2)}{\log_{5}x} = \log_{5}5 + \log_{5}x$\\

$\Leftrightarrow \log_{5}(6-x) + \log_{5}(x+2) = \log_{5}5x$\\

$\Leftrightarrow \log_{5}(6x+12-x^2-2x) = \log_{5}5x$\\

$\Leftrightarrow \log_{5}\left(\dfrac{-x^2+4x+12}{5x}\right) = 0$\\

$\Leftrightarrow \dfrac{-x^2+4x+12}{5x} = 1$\\

$\Leftrightarrow \dfrac{-x^2-x+12}{5x} = 0$

$\Delta = 49$\\

$x_{1,2} = \dfrac{1\pm7}{-2}$\\

$x_{1} = -4$ dit zit niet in het domein.\\

$x_{2} = 3$\\

Dus $x = 3$




\begin{thebibliography}{999}
%\bibitem{a} Referentie
\end{thebibliography}
\end{document} 
